% Authority: EGA I intro-fr.tex, SHA-256 CE78400CD9DBC36A4D11CC933B7BE18BEAC0C69AE6A04FBA3ECFA86053572980.
\section*{서론}

\begin{center}
\emph{오스카 자리스키와 앙드레 베유에게.}
\end{center}

\oldpage[I]{5}

이 논문과 뒤이어 나올 많은 논문은 대수기하학의 기초에 관한 하나의 논고를 이루기 위한 것이다.
원칙적으로 이 분야에 관한 어떤 특별한 지식도 전제하지 않는다. 오히려 그러한 지식은 분명한 장점이 있음에도 불구하고, 그것이 수반하는 쌍유리적 관점에 지나치게 익숙해진 탓에 여기서 제시하는 관점과 기법에 익숙해지려는 이에게 때로 방해가 될 수 있음이 드러났다.
반면 독자가 다음 주제들을 잘 알고 있다고 가정한다.

\begin{enumerate}
  \item[a)] \emph{가환대수학}. 예를 들어 현재 준비 중인 N.~부르바키의 \emph{원론} 여러 권에서 전개되는 내용이다(이 책들이 출간되기 전까지는 사뮈엘-자리스키~\cite{I-13}와 사뮈엘~\cite{I-11},~\cite{I-12}을 참조할 수 있다).
  \item[b)] \emph{호몰로지 대수학}. 카르탕-아일렌베르크~\cite{I-2}(이하 (M)), 고드망~\cite{I-4}(이하 (G)), 그리고 A.~그로텐디크의 최근 논문~\cite{I-6}(이하 (T))을 참조한다.
  \item[c)] \emph{층 이론}. 주된 참고문헌은 (G)와 (T)이다. 이 이론은 가환대수학의 핵심 개념들을 ``기하학적'' 용어로 해석하고 ``전역화''하는 데 필수적인 언어를 제공한다.
  \item[d)] 끝으로 독자가 \emph{함자 언어}에 어느 정도 익숙하면 유용할 것이다. 이 논고에서는 이 언어를 계속 사용한다. (M), (G), 특히 (T)를 참조할 수 있으며, 그 원리와 일반 함자론의 주요 결과들은 이 논고의 저자들이 준비 중인 별도의 저서에서 더 자세히 설명할 예정이다.
\end{enumerate}

\bigskip
\begin{center}*\quad*\quad*\end{center}
\bigskip

이 서론은 대수기하학에서 ``스킴''이라는 관점이 무엇인지 개략적으로 설명하거나, 이 관점을 채택해야 했던 수많은 이유를 열거할 자리가 아니다. 특히 우리가 다루는 ``다양체''의 국소환에서 멱영원을 체계적으로 받아들여야 하는 이유도 여기서는 설명하지 않는다. 이 수용은 필연적으로 유리 사상의 개념을 뒤로 물리고, 정칙 사상 또는 ``사상''의 개념을 앞세운다.
이 논고의 목적은 바로 ``스킴''의 언어를 체계적으로 전개하고, 바라건대 그 필요성을 입증하는 데 있다.
그렇게 하기는 어렵지 않겠지만, 제1장에서 전개하는 개념들에 대한 ``직관적'' 입문도 여기서는 시도하지 않겠다.

\oldpage[I]{6}
이 논고의 내용을 미리 훑어보고 싶은 독자는 A.~그로텐디크가 1958년 에든버러 국제수학자대회에서 한 강연~\cite{I-7}과 같은 저자의 발표~\cite{I-8}을 참조할 수 있다.
J.-P.~세르의 논문~\cite{I-14}(이하 (FAC))도 대수기하학의 고전적 관점과 스킴의 관점 사이를 잇는 설명으로 볼 수 있다. 따라서 이 논문을 읽는 것은 우리의 \emph{원론}을 읽기 위한 훌륭한 준비가 될 수 있다.

\bigskip
\begin{center}*\quad*\quad*\end{center}
\bigskip

참고로 이 논고에서 예정한 전체 구성을 아래에 제시한다. 특히 뒤쪽 장들의 구성은 나중에 바뀔 수 있다.

\medskip
\begin{tabular}{rl}
  제1장. & --- 스킴의 언어.\\
  --- \textbf{II}. & --- 몇몇 사상 부류에 대한 초보적 전역 연구.\\
  --- \textbf{III}. & --- 연접 대수적 층의 코호몰로지. 응용.\\
  --- \textbf{IV}. & --- 사상의 국소적 연구.\\
  --- \textbf{V}. & --- 스킴의 초보적 구성법.\\
  --- \textbf{VI}. & --- 강하 기법. 스킴 구성의 일반적 방법.\\
  --- \textbf{VII}. & --- 군 스킴, 주다발 공간.\\
  --- \textbf{VIII}. & --- 다발 공간의 미분론적 연구.\\
  --- \textbf{IX}. & --- 기본군.\\
  --- \textbf{X}. & --- 유수와 쌍대성.\\
  --- \textbf{XI}. & --- 교차이론, 천 특성류, 리만-로흐 정리.\\
  --- \textbf{XII}. & --- 아벨 스킴과 피카르 스킴.\\
  --- \textbf{XIII}. & --- 베유 코호몰로지.\\
\end{tabular}
\medskip

원칙적으로 모든 장은 열려 있는 것으로 간주하며, 언제든지 절을 덧붙일 수 있다. 채택한 출판 방식의 불편을 줄이기 위해 그러한 절은 별도의 분책으로 출간한다.
한 장이 출간될 때 어떤 절이 예정되어 있거나 준비 중이면, 긴급성에 따른 순서 때문에 실제 출간이 훨씬 뒤로 미뤄지더라도 그 장의 목차에 이를 적어 둔다.
독자의 편의를 위해 이 논고의 여러 장에서 사용하는 가환대수학, 호몰로지 대수학, 층 이론의 여러 보충 결과를 ``제0장''에 모았다. 이 결과들은 어느 정도 알려져 있지만 편리한 참고문헌을 제시할 수 없었던 것들이다.
독자는 논고 본문을 읽다가 우리가 인용하는 결과가 충분히 익숙하지 않을 때에만 제0장을 참조하기를 권한다.

\oldpage[I]{7}
이렇게 읽으면 이 논고는 초심자가 가환대수학과 호몰로지 대수학에 익숙해지는 좋은 방법이 될 수 있다고 생각한다. 구체적인 응용을 동반하지 않는 이 두 분야의 공부는 적지 않은 이들에게 따분하고 심지어 침울하기까지 한 것으로 여겨지기 때문이다.

\bigskip
\begin{center}*\quad*\quad*\end{center}
\bigskip

이 서론에서 여기 제시하는 개념과 결과의 역사를, 아무리 간략하게라도 개관하는 일은 우리의 능력 밖이다.
본문에는 이해에 특히 유용하다고 판단한 참고문헌만 싣고, 가장 중요한 결과들에 대해서만 그 기원을 밝힐 것이다.
적어도 형식적으로는 이 저서가 다루는 주제들이 상당히 새롭다. 그래서 우리가 그 연구를 전해 들었을 뿐인 19세기와 20세기 초 대수기하학의 선구자들을 거의 인용하지 않는 것이다.
그러나 저자들에게 가장 직접적인 영향을 주고 스킴 관점의 발전에 기여한 저작들에 관해서는 여기서 몇 마디 할 필요가 있다.
무엇보다 먼저 J.-P.~세르의 기초적인 논문 (FAC)을 들 수 있다. 이 논문은 A.~베유의 고전적인 \emph{Foundations}~\cite{I-18}가 너무 건조하다고 느낀 여러 젊은 입문자들(그 가운데에는 이 논고의 저자 한 명도 있다)에게 대수기하학의 입문서 역할을 했다.
여기서 처음으로 ``추상'' 대수다양체의 ``자리스키 위상''이 대수적 위상수학의 특정 기법들을 적용하기에 완전히 적합하며, 특히 코호몰로지 이론을 낳는다는 사실이 증명되었다.
더 나아가 그곳에서 주어진 대수다양체의 정의는 우리가 여기서 전개하는 개념의 확장에 가장 자연스럽게 맞는다\footnote{J.-P.~세르가 지적했듯이, 환의 층을 주어 다양체의 구조를 정의한다는 발상은 H.~카르탕에게서 비롯되었으며, 카르탕은 이를 해석공간 이론의 출발점으로 삼았다.
물론 대수기하학에서와 마찬가지로 ``해석기하학''에서도 해석공간의 국소환 안에 있는 멱영원에 정당한 지위를 부여해야 한다.
H.~카르탕과 J.-P.~세르의 정의를 이렇게 확장하는 문제는 최근 H.~그라우어트~\cite{I-5}가 다루기 시작했으며, 이 일반적 틀 안에서 해석기하학을 체계적으로 설명하는 저서가 머지않아 나오기를 기대할 수 있다.
이 논고에서 전개하는 개념과 기법이 해석기하학에서도 의미를 유지한다는 점은 분명하다. 다만 후자의 이론에서는 기술적 어려움이 훨씬 더 클 것으로 예상해야 한다.
방법이 단순한 대수기하학은 앞으로 해석공간 이론이 발전하는 데 일종의 형식적 모형이 될 수 있을 것이다.}.
또한 세르는 체 위의 아핀 대수 대신 임의의 가환환을 쓰면 아핀 대수다양체의 코호몰로지 이론을 어려움 없이 옮겨 적을 수 있음을 이미 알아차렸다.
따라서 이 논고의 제I장과 제II장, 그리고 제III장의 첫 두 절은 본질적으로 (FAC)와 같은 저자의 뒤이은 논문~\cite{I-15}의 주요 결과를 이 확장된 틀로 곧바로 옮긴 것으로 볼 수 있다.
우리는 C.~슈발레의 \emph{대수기하학 세미나}~\cite{I-1}에서도 큰 도움을 얻었다. 특히 그가 도입한 ``구성가능 집합''을 체계적으로 사용하는 것이 스킴 이론에서 매우 유용하다는 사실이 드러났다(제IV장 참조).
또한 차원의 관점에서 사상을 연구하는 방법(제IV장)도 그에게서 가져왔다. 이 방법은 스킴의 틀로 거의 아무 변화 없이 옮겨진다.

\oldpage[I]{8}
한편 슈발레가 도입한 ``국소환의 스킴''이라는 개념은 대수기하학을 자연스럽게 확장할 수 있게 한다. 다만 우리가 여기서 부여하려는 유연성과 일반성을 모두 갖추지는 못한다. 이 개념과 우리 이론의 관계는 제1장 \S~8을 보라.
M.~나가타는 데데킨트 환 위의 대수기하학에 관한 많은 특수한 결과를 담은 일련의 논문~\cite{I-9}에서 이러한 확장을 전개하였다\footnote{대수기하학에 관한 우리의 관점에 가까운 연구로는 E.~켈러의 중요한 저작~\cite{I-22}과 저우와 이구사의 최근 노트~\cite{I-3}을 들 수 있다. 이들은 나가타-슈발레 이론의 틀 안에서 (FAC)의 몇몇 결과를 다시 다루고 퀸네트 공식도 제시한다.}.

\bigskip
\begin{center}*\quad*\quad*\end{center}
\bigskip

끝으로 대수기하학에 관한 책, 특히 그 기초에 관한 책이 O.~자리스키와 A.~베유 같은 수학자들에게서, 설령 다른 사람들을 거쳐서라도, 필연적으로 영향을 받는다는 것은 말할 필요도 없다.
특히 자리스키의 \emph{정칙함수론}~\cite{I-20}은 코호몰로지 방법으로 적절히 유연하게 만들고 존재정리(제III장 \S\S~4, 5)로 보완하면, 제VI장에서 설명하는 강하 기법과 더불어 이 논고에서 사용하는 주요 도구 가운데 하나가 된다. 우리가 보기에 이는 대수기하학에서 쓸 수 있는 가장 강력한 도구 가운데 하나이다.

이 이론이 속하는 일반적 기법은 다음과 같이 개략적으로 설명할 수 있다. 전형적인 예는 제IX장에서 기본군을 연구할 때 나타난다.
한 대수다양체에서 다른 대수다양체로 가는 고유 사상(제II장) $f:X\to Y$가 있고(더 일반적으로는 한 스킴에서 다른 스킴으로 가는 사상), 점 $y\in Y$의 근방에 관한 문제 P를 풀기 위해 $y$의 근방에서 이 사상을 연구하려 한다.
다음 단계를 차례로 밟는다.
\begin{enumerate}
  \item[$1^\circ$] $Y$가 아핀이라고 가정할 수 있다. 그러면 $X$는 $Y$의 아핀환 $A$ 위에서 정의된 스킴이 되며, 나아가 $A$를 $y$의 국소환으로 바꿀 수 있다.
  이 축약은 실제로 언제나 쉽고(제V장), $A$가 \emph{국소환}인 경우로 문제를 돌린다.
  \item[$2^\circ$] $A$가 \emph{아르틴 국소환}인 경우에 문제를 연구한다.
  $A$가 정역이라고 가정하지 않을 때에도 문제가 실제로 의미를 유지하게 하려면 문제 P를 다시 정식화해야 할 때가 있다. 이렇게 하면 흔히 이 단계에서 ``무한소적'' 성격을 띠는 문제를 더 잘 이해하게 된다.
  \item[$3^\circ$] 형식 스킴 이론(제III장 \S\S~3, 4, 5)을 이용하면 아르틴 환의 경우에서 \emph{완비 국소환}의 경우로 넘어갈 수 있다.
  \item[$4^\circ$] 끝으로 $A$가 임의의 국소환이면, $X$ 위의 적당한 스킴에서 주어진 ``형식적'' 절단에 가까운 ``다중 절단''을 고려함으로써(제IV장), $A$의 완비화로 스칼라를 확대하여 $X$에서 얻은 스킴에 대해 알려진 결과를 $A$의 충분히 단순한 유한 확대(예를 들면 비분기 확대)에 대한 유사한 결과로 흔히 옮길 수 있다.
\end{enumerate}

이 개요는 아르틴 환 $A$ 위에서 정의된 스킴을 체계적으로 연구하는 일이 중요함을 보여 준다.
세르가 국소 유체론을 정식화할 때 취한 관점과 그린버그의 최근 연구는, 그러한 스킴 $X$에 $A$의 잉여체 $k$(완전체라고 가정한다) 위의 스킴 $X'$를 함자적으로 대응시키는 방식으로 이 연구를 수행할 수 있음을 시사하는 듯하다. 여기서 $X'$의 차원은 좋은 경우 $n\dim X$이고, $n$은 $A$의 길이이다.

A.~베유의 영향에 관해서는 다음과 같이 말하는 것으로 충분하다. ``베유 코호몰로지''의 정의를 필요한 일반성 전체를 갖추어 정식화하고, 디오판토스 기하학의 유명한 추측들~\cite{I-19}을 확립하는 데 필요한 모든 형식적 성질의 증명\footnote{오해를 피하기 위해 덧붙이면, 이 서론을 쓰는 시점에는 이 과제가 이제 막 시작되었을 뿐이며, 따라서 아직 베유 추측의 증명에 이르지 못했다.}에 착수하려면 필요한 도구를 개발해야 한다. 이것이 이 논고를 집필하게 된 주된 동기 가운데 하나였다. 대수기하학에서 통상 사용하는 개념과 방법의 자연스러운 틀을 찾고, 저자들 자신이 그 개념과 기법을 이해할 기회를 얻으려는 바람도 똑같이 중요한 동기였다.

\bigskip
\begin{center}*\quad*\quad*\end{center}
\bigskip

끝으로 독자들에게 미리 알려 두는 것이 유익하다고 생각한다. 저자들 자신이 그랬듯이, 독자들도 스킴의 언어에 익숙해지고 기하학적 직관이 제안하는 통상적 구성들을 본질적으로 단 하나의 합리적인 방식으로 이 언어에 옮길 수 있음을 확신하기까지는 어느 정도 어려움을 겪을 것이다.
현대수학의 많은 분야에서 그렇듯이, 최초의 직관은 그것을 필요한 정밀성과 일반성을 모두 갖추어 표현하는 고유한 언어에서 겉보기에는 점점 더 멀어지고 있다.
여기서 심리적 어려움은 집합에 익숙한 개념들, 곧 데카르트 곱, 군·환·가군의 연산, 다발, 주동차다발 등을 집합의 범주와 이미 상당히 다른 범주의 대상들(즉 준스킴의 범주 또는 주어진 준스킴 위의 준스킴 범주)로 옮겨야 한다는 데 있다.
미래의 수학자가 이러한 새로운 추상화의 노력을 피하기는 어려울 것이다. 그러나 집합론에 익숙해지면서 우리의 선배들이 기울였던 노력과 비교하면, 결국 이 노력은 꽤 작은 것일지도 모른다.

\bigskip
\begin{center}*\quad*\quad*\end{center}
\bigskip

참조는 십진법 체계에 따라 표시한다. 예를 들어 III, 4.9.3에서 III은 장, 4는 절, 9는 그 절의 항을 나타낸다.
같은 장 안에서는 장 표시를 생략한다.
